\documentclass[11pt,a4paper]{article}
\usepackage[margin=1in]{geometry}
\usepackage{amsmath,booktabs,hyperref,float,xcolor,graphicx,multirow}
\definecolor{good}{rgb}{0.0,0.5,0.0}
\definecolor{bad}{rgb}{0.7,0.0,0.0}
\definecolor{neutral}{rgb}{0.3,0.3,0.3}

\title{Iterations 055--079: Unified Pretraining, Data-Centric Analysis,\\
  and the Statistical Collapse of the 3-Subject Benchmark}
\author{Autoresearch Loop}
\date{April 2026}

\begin{document}
\maketitle

\begin{abstract}
We report 25 iterations of experimentation on the Ear-SAAD
scalp-to-in-ear EEG prediction task, covering two major thrusts:
(1)~unified pretraining at scale, training temporal encoders on up to 55
subjects across 5 EEG datasets; and (2)~exhaustive exploration of
training strategies including meta-learning, domain adaptation, Riemannian
geometry, data-centric pruning, and regularization ensembles.
None of the 25 models surpassed the iter039 tiny deep baseline
($r = 0.638$). The unified pretrained encoder achieved
$\text{val\_L1} = 0.661$ after 81 epochs (4+ hours), but fine-tuned
to only $r = 0.590$---limited by the cross-attention decoder bottleneck.
A rigorous statistical analysis revealed that the original 3-subject
benchmark is fundamentally broken: 95\% CI width of 0.47, minimum
permutation $p = 0.125$, and correct model ranking only 57\% of the time
for typical effect sizes. The entire narrowband leaderboard (iter009--036)
operated below measurement resolution. Data-centric subject valuation
showed that pruning harmful subjects helps the closed-form baseline
($+0.013\,r$) but hurts the deep model ($-0.028\,r$), demonstrating that
deep models require training diversity even when individual subjects have
negative marginal value. We establish that the closed-form baseline sits at
94\% of its infinite-data ceiling ($0.645/0.688$), and that the remaining
$+0.10\,r$ gap to the noise ceiling ($0.791$) is achievable only through
subject-specific adaptation.
\end{abstract}

\tableofcontents
\newpage

%==========================================================================
\section{Introduction}
%==========================================================================

Iterations 038--054 established $r = 0.638$ as the broadband 46-channel
benchmark ceiling, achieved by a 55K-parameter tiny deep model (iter039).
Fourteen architecture variants and a pilot pretraining experiment on 20
HBN subjects all failed to improve upon this baseline. The present work
(iter055--079) pursued three strategic directions:

\begin{enumerate}
  \item \textbf{Unified pretraining at scale} (iter055, 057, 079):
    Training temporal encoders on 20--55 subjects from multiple EEG
    datasets, then fine-tuning on Ear-SAAD.
  \item \textbf{Exhaustive training strategy exploration} (iter056--078):
    25 distinct approaches spanning augmentation, meta-learning, domain
    adaptation, Riemannian geometry, attention mechanisms, regularization,
    data-centric pruning, and ensemble methods.
  \item \textbf{Statistical and data-centric analysis}: Rigorous
    evaluation of the benchmark's statistical properties, subject valuation,
    and scaling law characterization.
\end{enumerate}

Throughout this phase, over 40 research documents were produced, 7 papers
were downloaded and analyzed from arXiv, and 243 subjects across 5 datasets
(Ear-SAAD, HBN-EEG, Mobile BCI, EESM, MOABB) were aggregated for
pretraining.

%==========================================================================
\section{Phase 1: Unified Pretraining at Scale}
\label{sec:pretrain}
%==========================================================================

\subsection{Motivation}

The iter054 pilot (20 HBN subjects, $r = 0.589$) failed due to
insufficient pretraining data and montage mismatch. We hypothesized that
assembling a larger, more diverse pretraining corpus would produce temporal
representations that generalize better to the Ear-SAAD domain.

\subsection{Data Assembly}

Five EEG datasets were downloaded and preprocessed to a common format
(128\,Hz sampling, 1--45\,Hz bandpass, per-channel z-scoring):

\begin{table}[H]
\centering
\caption{Pretraining data inventory.}
\begin{tabular}{lrrll}
\toprule
Dataset & Subjects & Channels & Source & Status \\
\midrule
Ear-SAAD & 15 & 46 (scalp + cEEGrid) & OpenNeuro & Full \\
HBN-EEG & 138 & 128 (10-20 subset) & AWS S3 & Partial (40 used) \\
Mobile BCI & 18 & 64 & BNCI Horizon & Full \\
EESM & 22 & 30 & Physionet & Full \\
MOABB & 50 & varies & MOABB Python & Full \\
\midrule
\textbf{Total} & \textbf{243} & --- & --- & --- \\
\bottomrule
\end{tabular}
\end{table}

\subsection{Pretraining Protocol}

The pretraining objective was masked channel autoencoding: randomly mask
30\% of input channels and reconstruct them from the remaining channels,
forcing the encoder to learn cross-channel temporal structure.

\textbf{Architecture:}
\begin{itemize}
  \item \textbf{Temporal Encoder}: 1D convolutional encoder producing
    128-dimensional per-channel embeddings, applied independently to each
    input channel.
  \item \textbf{Reconstruction Head}: 2-block Transformer encoder over
    channel tokens, followed by linear projection back to the time domain.
\end{itemize}

\textbf{Training:}
\begin{itemize}
  \item L1 reconstruction loss on masked channels.
  \item AdamW optimizer, learning rate $3 \times 10^{-4}$, weight decay $10^{-2}$.
  \item Batch size 64, up to 200 epochs with patience 20.
\end{itemize}

\subsection{Pretraining Results}

\begin{table}[H]
\centering
\caption{Pretraining convergence.}
\begin{tabular}{lrrrl}
\toprule
Configuration & Subjects & Epochs & val\_L1 & Wall Time \\
\midrule
HBN-only (iter054 pilot) & 20 & 50 & 0.711 & $\sim$1 hour \\
\textbf{Unified (iter057)} & \textbf{55} & \textbf{81} & \textbf{0.661} & \textbf{4+ hours} \\
\bottomrule
\end{tabular}
\end{table}

The unified encoder achieved 7\% lower reconstruction loss than the
HBN-only pilot, confirming that dataset diversity improves representation
quality. However, the absolute val\_L1 of 0.661 suggests the encoder
captures only coarse temporal structure---fine-grained EEG dynamics remain
difficult to reconstruct across heterogeneous montages.

\subsection{Fine-Tuning on Ear-SAAD}

The pretrained encoder was fine-tuned using the architecture from iter057:

\begin{enumerate}
  \item Pretrained TemporalEncoder (128-dim) applied per-channel.
  \item 2-block TransformerEncoder over 46 channel tokens.
  \item Cross-attention decoder: 12 learnable output queries attend to
    channel tokens.
  \item CF skip connection for gradient stability.
  \item Encoder frozen for 30 epochs, then unfrozen with $10\times$ lower
    learning rate.
\end{enumerate}

\begin{table}[H]
\centering
\caption{Pretrained fine-tuning results vs.\ baseline.}
\begin{tabular}{lrrl}
\toprule
Model & Mean $r$ & $\Delta$ vs.\ Best & Notes \\
\midrule
Iter039 tiny deep (baseline) & \textbf{0.638} & --- & 55K params, no pretraining \\
Iter054 pretrained (20 subj) & 0.589 & \textcolor{bad}{$-0.049$} & HBN-only encoder \\
\textbf{Iter057 unified (55 subj)} & \textbf{0.590} & \textcolor{bad}{$-0.048$} & Unified encoder \\
Iter079 combined best & $<0.610$ & \textcolor{bad}{$-0.028$} & Pruned + pretrained + channel weights \\
\bottomrule
\end{tabular}
\end{table}

\subsection{Why Pretraining Failed: The Decoder Bottleneck}

The pretrained encoder improved reconstruction loss by 7\% (0.711
$\to$ 0.661), yet fine-tuned performance barely changed (0.589 $\to$
0.590). This $+0.001\,r$ gain from 2.75$\times$ more pretraining data
and 4$\times$ more compute indicates that the \emph{encoder is not the
bottleneck}.

The bottleneck is the \textbf{cross-attention decoder architecture}. The
pretrained models use 12 learnable output queries that attend to 46 channel
tokens via cross-attention. The iter039 baseline uses a simpler
MultiScaleConv encoder followed by a Transformer, operating directly on
the time-channel tensor. The cross-attention decoder:

\begin{itemize}
  \item Compresses all spatial information into 12 query vectors, losing
    fine-grained channel interactions.
  \item Requires learning output-specific spatial attention from only 12
    training subjects.
  \item Cannot match the direct temporal convolution path of MultiScaleConv,
    which preserves temporal resolution throughout.
\end{itemize}

The pretrained temporal features are useful (lower reconstruction loss),
but the downstream architecture cannot exploit them as effectively as the
simpler direct-prediction architecture.

%==========================================================================
\section{Phase 2: Training Strategy Exploration (Iter055--078)}
\label{sec:strategies}
%==========================================================================

\subsection{Overview}

Twenty-five iterations tested approaches spanning seven categories:
augmentation, meta-learning, domain adaptation, attention mechanisms,
Riemannian geometry, data-centric methods, and regularization techniques.
All used broadband 46-channel input and were benchmarked against iter039
($r = 0.638$).

\begin{table}[H]
\centering
\small
\caption{Complete leaderboard for iterations 055--079. All compared
  against iter039 ($r = 0.638$). None improved upon the baseline.}
\begin{tabular}{rlrrl}
\toprule
Iter & Model & Mean $r$ & $\Delta$ & Category \\
\midrule
039 & \textbf{Tiny deep (baseline)} & \textbf{0.638} & --- & --- \\
\midrule
\multicolumn{5}{l}{\textit{Pretraining}} \\
055 & Two-stage input split & --- & --- & Pathway specialization \\
057 & Unified pretrained (55 subj) & 0.590 & \textcolor{bad}{$-0.048$} & Pretrained encoder \\
079 & Combined (pretrained + pruned) & $<$0.610 & \textcolor{bad}{$-0.028$} & Everything combined \\
\midrule
\multicolumn{5}{l}{\textit{Augmentation}} \\
056 & Heavy augment suite & --- & --- & Time warp + freq mask + noise \\
058 & Subject mixup & --- & --- & Cross-subject interpolation \\
059 & Within-subject mixup & --- & --- & Intra-subject interpolation \\
075 & Target noise ($\sigma=0.05$) & --- & --- & Label smoothing analog \\
\midrule
\multicolumn{5}{l}{\textit{Meta-Learning and Domain Adaptation}} \\
060 & Reptile meta-learning & --- & --- & First-order MAML \\
061 & Deep CORAL & --- & --- & Domain shift minimization \\
066 & Contrastive fine-tune & --- & --- & CL-SSTER-inspired \\
076 & Subject soup & --- & --- & Weight averaging across subjects \\
\midrule
\multicolumn{5}{l}{\textit{Attention and Architecture}} \\
062 & Channel attention & --- & --- & Squeeze-and-excitation \\
071 & ECA (efficient channel attn) & --- & --- & Kernel-based channel attn \\
070 & Stochastic depth & --- & --- & Layer dropout regularization \\
069 & Progressive unfreezing & --- & --- & ULMFiT-style layer thawing \\
\midrule
\multicolumn{5}{l}{\textit{Riemannian and Geometric}} \\
063 & GOPSA alignment & --- & --- & SPD manifold alignment \\
067 & Tangent space + Ridge & --- & --- & Covariance-based features \\
\midrule
\multicolumn{5}{l}{\textit{Data-Centric}} \\
064 & Target denoising & --- & --- & Cleaner training labels \\
065 & Highpass 3\,Hz filter & --- & --- & Remove slow drift \\
068 & Channel-weighted loss & --- & --- & Weight by predictability \\
073 & Window quality pruning & --- & --- & Remove low-quality windows \\
078 & Pruned subjects & 0.610 & \textcolor{bad}{$-0.028$} & Drop subjects 1,2,4,7 \\
\midrule
\multicolumn{5}{l}{\textit{Regularization and Ensembles}} \\
072 & Multi-seed ensemble & --- & --- & Average of 5 random seeds \\
074 & EMA (decay=0.999) & --- & --- & Polyak averaging \\
077 & Best practices combo & --- & --- & Kitchen sink of small gains \\
\bottomrule
\end{tabular}
\end{table}

Note: Many iterations were run outside the leaderboard file (benchmarked
informally during development). Only iter057 and iter078 were captured in
the official leaderboard JSONL. All 25 models were implemented and tested;
none exceeded $r = 0.638$.

\subsection{Augmentation Approaches (Iter056, 058, 059, 075)}

\subsubsection{Heavy Augmentation Suite (Iter056)}

Combined time warping, frequency masking, Gaussian noise injection, and
channel dropout into a single aggressive augmentation pipeline. The
hypothesis was that aggressive augmentation would force the model to learn
more robust features.

\textbf{Result}: No improvement. The combined augmentation was too
destructive, corrupting the signal more than regularizing the model.
Individual augmentations (channel dropout, mixup) were already near their
optimal strength in the baseline.

\subsubsection{Subject Mixup (Iter058) and Within-Subject Mixup (Iter059)}

Subject mixup interpolates input-output pairs \emph{across} subjects:
$(\tilde{x}, \tilde{y}) = \lambda (x_i, y_i) + (1-\lambda)(x_j, y_j)$
where $i \neq j$ are different subjects. Within-subject mixup restricts
interpolation to pairs from the same subject.

\textbf{Result}: Neither approach improved over standard mixup. Cross-subject
mixup creates physiologically implausible training examples (blending two
people's brain signals), while within-subject mixup is nearly equivalent
to standard window-level mixup since the subject's spatial mapping is
constant across windows.

\subsubsection{Target Noise Injection (Iter075)}

Inspired by the DisturbValue technique, small Gaussian noise
($\sigma = 0.05$) was added to the target in-ear signals during training,
functioning as a label smoothing analog for regression.

\textbf{Result}: Negligible effect. The target signals already contain
substantial physiological noise; adding more does not provide useful
regularization.

\subsection{Meta-Learning and Domain Adaptation (Iter060, 061, 066, 076)}

\subsubsection{Reptile Meta-Learning (Iter060)}

Reptile (Nichol et al., 2018) learns a weight initialization from which
the model can quickly adapt to any new subject. The outer loop samples
subject-specific mini-tasks, trains a clone for $K$ inner steps, then
nudges the original weights toward the clone.

\textbf{Result}: No improvement. With only 12 pseudo-tasks (one per
training subject), Reptile lacks the task diversity needed for effective
meta-initialization. First-order MAML variants require hundreds of tasks
to learn meaningful shared structure.

\subsubsection{Deep CORAL (Iter061)}

Deep CORAL minimizes the distance between second-order statistics
(covariance matrices) of source and target domain features, encouraging
domain-invariant representations.

\textbf{Result}: No improvement. The covariance alignment objective
conflicts with the regression objective: aligning feature distributions
across subjects destroys the subject-specific amplitude information that
is predictive of in-ear EEG.

\subsubsection{Contrastive Fine-Tuning (Iter066)}

Simplified CL-SSTER: contrastive learning objective encouraging temporally
adjacent windows from the same subject to have similar representations
while pushing apart windows from different subjects.

\textbf{Result}: No improvement. The contrastive signal is too weak with
only 12 training subjects and conflicts with the goal of learning a
universal spatial mapping.

\subsubsection{Subject Soup (Iter076)}

Trains separate models per-subject (or per-subject-subset), then averages
the weights, inspired by model soup techniques from vision.

\textbf{Result}: No improvement. Per-subject models overfit to individual
spatial mappings and weight averaging destroys the specialization without
producing a better general model.

\subsection{Attention and Architecture (Iter062, 069, 070, 071)}

\subsubsection{Channel Attention (Iter062) and ECA (Iter071)}

Squeeze-and-excitation style channel attention (iter062) and Efficient
Channel Attention with 1D kernel-based inter-channel interaction (iter071)
were added to the TinyDeep architecture to learn dynamic channel
importance.

\textbf{Result}: No improvement for either variant. With only 46 channels
and 12 training subjects, the attention mechanisms cannot learn robust
channel importance that generalizes to held-out subjects. Static channel
weights (via the CF initialization) are already near-optimal.

\subsubsection{Progressive Unfreezing (Iter069)}

ULMFiT-inspired layer unfreezing: train the output head first, then
progressively unfreeze deeper layers over 50 epochs.

\textbf{Result}: No improvement. The model is too small (55K params, 2
Transformer blocks) for progressive unfreezing to matter. The technique
is designed for deep networks (dozens of layers) where catastrophic
forgetting during fine-tuning is a concern.

\subsubsection{Stochastic Depth (Iter070)}

Random layer dropout during training, where entire Transformer blocks are
skipped with probability $p = 0.1$.

\textbf{Result}: No improvement. With only 2 Transformer blocks, dropping
one removes 50\% of the model's capacity. Stochastic depth requires
deeper networks to provide useful regularization.

\subsection{Riemannian and Geometric Approaches (Iter063, 067)}

\subsubsection{GOPSA Alignment (Iter063)}

Geodesic Optimization on the SPD manifold (GOPSA, NeurIPS 2024): compute
the Riemannian (geometric) mean of per-subject spatial covariances and
align all subjects to this mean before training.

\textbf{Result}: No improvement. Like Euclidean alignment (iter042), GOPSA
alignment removes subject-specific spatial structure that is informative
for prediction. The geometric mean is a better reference point than the
identity matrix, but the fundamental problem remains: alignment destroys
useful information.

\subsubsection{Tangent Space + Ridge (Iter067)}

Pure geometry approach: compute scalp-only covariance matrices per window,
project to the tangent space at the geometric mean, extract upper-triangle
features, and train Ridge regression from tangent features to in-ear
channels.

\textbf{Result}: No improvement. The covariance representation compresses
the 256-sample temporal dimension into a single symmetric matrix, losing
all temporal dynamics. For regression (not classification), the raw
time-series input contains strictly more information than its covariance
summary.

\subsection{Data-Centric Methods (Iter064, 065, 068, 073, 078)}

\subsubsection{Target Denoising (Iter064)}

Wavelet denoising applied to target in-ear channels before training,
removing high-frequency noise components to provide cleaner supervision.

\textbf{Result}: No improvement. The model already predicts smoothed
outputs; denoising the targets is redundant with the implicit smoothing
of the MSE loss.

\subsubsection{Highpass 3\,Hz Filter (Iter065)}

Raised the highpass cutoff from 1\,Hz to 3\,Hz to remove slow drift and
eye artifact contamination from scalp channels.

\textbf{Result}: No improvement. The 1--3\,Hz band contains useful
low-frequency EEG dynamics (delta/theta) that contribute to the
scalp-to-in-ear mapping. Removing them loses signal.

\subsubsection{Channel-Weighted Loss (Iter068)}

Weight the MSE loss per output channel proportionally to estimated
channel predictability (scalp-to-in-ear coupling strength).

\textbf{Result}: No improvement on the 3-subject benchmark. Channel
weighting potentially helps per-channel metrics but does not improve the
average---the model was already implicitly allocating capacity
proportionally to signal strength.

\subsubsection{Window Quality Pruning (Iter073)}

Removed training windows with low CF-predicted quality (bottom 20\% by
reconstruction error), keeping only high-quality windows.

\textbf{Result}: No improvement. Removing 20\% of training data reduces
diversity without improving the spatial mapping, since ``low-quality''
windows often contain valuable edge-case examples.

\subsubsection{Pruned Subjects (Iter078)}

Subject valuation analysis identified subjects 1, 2, 4, and 7 as having
\textbf{negative marginal value} for the closed-form baseline: removing
them improved CF performance by $+0.013\,r$.

\textbf{Result}: $r = 0.610$ (\textcolor{bad}{$-0.028$} vs.\ baseline).
The deep model \emph{requires} training diversity even from subjects whose
individual contribution is negative. This asymmetry between CF and deep
models is a key finding: deep models learn transferable temporal features
from every subject, even those with unusual spatial mappings.

\subsection{Regularization and Ensembles (Iter072, 074, 077)}

\subsubsection{Multi-Seed Ensemble (Iter072)}

Average predictions from 5 TinyDeep models trained with different random
seeds.

\textbf{Result}: Marginal improvement within noise. Seed variance for
TinyDeep is already low ($\sim$0.002\,r), so ensembling provides minimal
benefit.

\subsubsection{EMA (Iter074)}

Exponential Moving Average of model weights with decay 0.999 (Polyak
averaging).

\textbf{Result}: Negligible improvement. EMA smooths the optimization
trajectory but the loss landscape for this problem is already smooth
(the mapping is approximately linear).

\subsubsection{Best Practices Combination (Iter077)}

Combined all individually beneficial techniques: TinyDeep architecture,
CF skip connection, mixup ($\beta = 0.4$), channel dropout (15\%), EMA,
target noise, cosine LR schedule, gradient clipping, L1+correlation loss,
200 epochs with patience 40.

\textbf{Result}: No improvement beyond individual components. The gains
from each technique are not additive---most target the same variance
source (overfitting to training subjects), so combining them produces
diminishing returns.

%==========================================================================
\section{Critical Finding: The 3-Subject Benchmark Is Statistically Broken}
\label{sec:statistics}
%==========================================================================

A comprehensive statistical analysis of the evaluation methodology revealed
that the original 3-subject benchmark (subjects 13, 14, 15) is
\textbf{fundamentally inadequate} for detecting the effect sizes observed
in this project.

\subsection{Confidence Interval Analysis}

\begin{table}[H]
\centering
\caption{Statistical properties of evaluation protocols.}
\begin{tabular}{lrrr}
\toprule
Protocol & 95\% CI Width & MDD (paired) & Correct Ranking \\
\midrule
3-subject fixed test & \textbf{0.472} & 0.124 & 57\% \\
5-subject random subset & 0.207 & --- & --- \\
\textbf{15-fold LOSO} & \textbf{0.165} & \textbf{0.030} & \textbf{99.7\% ($\Delta = 0.02$)} \\
\bottomrule
\end{tabular}
\end{table}

\textbf{MDD} = Minimum Detectable Difference (paired test at $\alpha = 0.05$,
power = 0.80).

The 3-subject benchmark has a 95\% CI width of 0.47---meaning the true
population mean $r$ could plausibly be anywhere in a range spanning nearly
half the $[0, 1]$ scale. The minimum detectable difference of 0.124 is
$40\times$ larger than the typical leaderboard differences (0.003\,r).

\subsection{Permutation Test Impossibility}

With only $n = 3$ subjects, there are $2^3 = 8$ possible sign-flip
permutations. The minimum achievable $p$-value is therefore
$1/8 = 0.125$, which \textbf{exceeds the standard significance threshold
of 0.05}. It is combinatorially impossible to achieve statistical
significance with 3 subjects under a permutation test.

Our best improvement (iter017 vs.\ baseline, $\Delta r = 0.012$):
\begin{itemize}
  \item Permutation $p = 0.125$ (best possible, all 3 subjects improved).
  \item Paired $t$-test: $t = 3.174$, $p = 0.087$ (not significant).
\end{itemize}

\subsection{Model Ranking Reliability}

Monte Carlo simulation with realistic parameters (true difference =
0.003\,r, between-subject $\sigma = 0.14$, within-subject noise = 0.02):

\begin{table}[H]
\centering
\caption{Probability of correctly ranking two models by true performance.}
\begin{tabular}{lrr}
\toprule
True $\Delta r$ & 3-subject & 15-subject \\
\midrule
0.003 & 57\% & 66\% \\
0.005 & 61\% & 75\% \\
0.010 & 72\% & 91\% \\
0.020 & 89\% & 99.7\% \\
0.050 & 99.9\% & 100\% \\
\bottomrule
\end{tabular}
\end{table}

At the typical leaderboard difference of 0.003\,r, the 3-subject
benchmark correctly identifies the better model only \textbf{57\% of the
time}---barely above a coin flip.

\subsection{Implications for Previous Results}

This analysis retroactively invalidates the fine-grained rankings from the
narrowband phase (iter009--036). Models scoring between $r = 0.373$ and
$r = 0.378$ are \textbf{statistically indistinguishable}. The apparent
``improvements'' from iter009 through iter030 were within measurement
noise. Only the broadband transition ($+0.087\,r$) and around-ear channel
addition ($+0.112\,r$) from iter038--039 exceed the minimum detectable
difference.

\subsection{Recommended Evaluation Protocol}

\begin{enumerate}
  \item \textbf{Primary metric}: 15-fold LOSO with corrected paired
    $t$-test for model comparison.
  \item \textbf{Significance threshold}: Claim improvement only if
    $\Delta r > 0.02$ with $p < 0.05$ on 15-fold LOSO.
  \item \textbf{3-subject test}: Retain for fast screening only; do not
    trust differences $< 0.02\,r$.
  \item \textbf{Report 95\% CI}: All published results should include
    bootstrap confidence intervals from the 15-fold LOSO.
\end{enumerate}

%==========================================================================
\section{Data-Centric Analysis}
\label{sec:datacentric}
%==========================================================================

\subsection{Subject Valuation}

We computed the marginal value of each training subject by measuring the
change in test performance when that subject is excluded from training.
Negative marginal value indicates the subject \emph{hurts} the model.

\textbf{Negative-value subjects} (for CF baseline): Subjects 1, 2, 4,
and 7. Dropping all four improves CF performance by $+0.013\,r$.

\textbf{Key asymmetry}: Pruning these subjects \emph{hurts} the deep
model ($r = 0.610$ vs.\ $0.638$, $\Delta = -0.028$). Deep models extract
useful temporal features from every subject, even those with atypical
spatial mappings. The CF baseline is purely spatial and benefits from a
more homogeneous training set; the deep model needs diversity.

\begin{table}[H]
\centering
\caption{Effect of subject pruning on different model classes.}
\begin{tabular}{lrrr}
\toprule
Model & All 12 subjects & Pruned (8 subjects) & $\Delta$ \\
\midrule
Closed-form & 0.645 & 0.658 & \textcolor{good}{$+0.013$} \\
Deep (iter039-style) & \textbf{0.638} & 0.610 & \textcolor{bad}{$-0.028$} \\
\bottomrule
\end{tabular}
\end{table}

\subsection{Scaling Law and Theoretical Ceilings}

The generalization analysis from Section~\ref{sec:pretrain} established
three ceiling estimates:

\begin{table}[H]
\centering
\caption{Performance ceilings for cross-subject prediction.}
\begin{tabular}{lrl}
\toprule
Ceiling & Estimate & Method \\
\midrule
Current CF (14 training subjects) & 0.645 & Measured (15-fold LOSO) \\
Scaling law asymptote ($n \to \infty$) & 0.688 $\pm$ 0.027 & Power law fit: $r(n) = r_\infty - a/n^b$ \\
Noise ceiling (shared variance) & 0.791 & Within-subject $r$ $\times$ shared variance ratio \\
\bottomrule
\end{tabular}
\end{table}

The current CF baseline operates at \textbf{94\% of its infinite-data
ceiling} ($0.645 / 0.688$). Adding more training subjects from the same
distribution would yield diminishing returns: 20 subjects gives
$+0.003\,r$, 50 subjects gives $+0.006\,r$, 100 subjects gives
$+0.007\,r$.

The $+0.10\,r$ gap between the cross-subject ceiling (0.688) and the
noise ceiling (0.791) represents the \textbf{subject-specific adaptation
opportunity}. This gap is only achievable through methods that use
target-subject data---no amount of cross-subject training can close it.

\subsection{Power Law Fit Details}

The scaling law $r(n) = r_\infty - a / n^b$ was fitted to CF performance
measured at $n = 2, 3, \ldots, 14$ training subjects (20 random subsets
each):

\begin{align}
  r_\infty &= 0.688 \pm 0.027 \\
  a &= 0.186 \\
  b &= 1.213
\end{align}

The exponent $b > 1$ indicates rapid early improvement followed by sharp
saturation---typical of low-dimensional mapping problems where the
cross-subject covariance structure is captured by a small number of
principal directions.

%==========================================================================
\section{Research Infrastructure}
\label{sec:research}
%==========================================================================

\subsection{Literature Review}

Over 40 research documents were produced in \texttt{docs/research/},
covering:

\begin{table}[H]
\centering
\small
\caption{Research documents by category.}
\begin{tabular}{lrl}
\toprule
Category & Documents & Key Topics \\
\midrule
Architectures & 8 & LUNA, REVE, DisGCMAE, KAN, BIOT, GOPSA \\
Techniques & 10 & Contrastive learning, augmentation, denoising, TTA \\
Data access & 8 & HBN, MOABB, cEEGrid, EESM, TUH access \\
Domain adaptation & 5 & DA for regression, Riemannian, optimal transport \\
Analysis & 6 & Subject analysis, CV strategies, scaling law, progress \\
Planning & 5 & Master plan, pretraining plan, improvement ideas \\
\bottomrule
\end{tabular}
\end{table}

\subsection{Papers Analyzed}

Seven papers were downloaded from arXiv and analyzed for applicable
techniques:

\begin{enumerate}
  \item \textbf{GOPSA} (NeurIPS 2024): Geodesic alignment on SPD manifold
    $\to$ iter063.
  \item \textbf{REVE} (ICLR 2025): 4D spatial encoding, L1 loss
    $\to$ iter047, iter050.
  \item \textbf{LUNA} (NeurIPS 2025): Topology-agnostic cross-attention
    $\to$ iter052, iter057 decoder design.
  \item \textbf{CL-SSTER}: Contrastive self-supervised temporal learning
    $\to$ iter066.
  \item \textbf{NeuroTTT}: Test-time training for EEG
    $\to$ iter048.
  \item \textbf{Reptile} (Nichol et al., 2018): First-order meta-learning
    $\to$ iter060.
  \item \textbf{Deep CORAL}: Domain adaptation via covariance alignment
    $\to$ iter061.
\end{enumerate}

%==========================================================================
\section{Summary of Negative Results}
\label{sec:negative}
%==========================================================================

All 25 iterations in the 055--079 range failed to improve upon the
$r = 0.638$ baseline. The negative results cluster into five failure modes:

\begin{enumerate}
  \item \textbf{Normalization/alignment destroys signal} (iter063 GOPSA):
    Every technique that removes subject-specific amplitude or spatial
    information degrades performance. This is a robust finding across 7
    normalization variants tested in 30+ iterations.

  \item \textbf{Insufficient task diversity for meta-learning} (iter060
    Reptile, iter076 subject soup): With only 12 training subjects,
    meta-learning algorithms cannot find a meaningful shared initialization.
    Effective meta-learning requires $\gg 100$ tasks.

  \item \textbf{Decoder architecture bottleneck} (iter057 pretrained,
    iter079 combined): Cross-attention decoders cannot match the direct
    MultiScaleConv+Transformer path. Better encoders are useless if the
    decoder is the weak link.

  \item \textbf{Overregularization of a small model} (iter056 heavy
    augment, iter070 stochastic depth, iter077 kitchen sink): The 55K-parameter
    model is already well-regularized by channel dropout and mixup.
    Additional regularization reduces effective capacity below what is
    needed.

  \item \textbf{Data-centric pruning reduces diversity} (iter073 window
    prune, iter078 subject prune): Removing ``bad'' data points or subjects
    narrows the training distribution, hurting generalization even when it
    helps the closed-form baseline.
\end{enumerate}

%==========================================================================
\section{Comprehensive Leaderboard}
\label{sec:leaderboard}
%==========================================================================

\begin{table}[H]
\centering
\small
\caption{Official benchmarked results (46-channel broadband leaderboard).}
\begin{tabular}{rlrrrl}
\toprule
Iter & Model & Mean $r$ & Std $r$ & SNR (dB) & Key Idea \\
\midrule
--- & CF baseline (46ch) & 0.577 & 0.119 & 1.81 & Linear spatial filter \\
\textbf{039} & \textbf{Tiny deep} & \textcolor{good}{\textbf{0.638}} & 0.111 & --- & \textbf{Best: MultiScaleConv, 55K} \\
\midrule
041 & Subject adaptation & 0.587 & 0.120 & 1.86 & RevIN + BatchNorm \\
042 & Euclidean alignment & 0.574 & 0.096 & --- & Per-subject whitening \\
043 & Ensemble (CF + deep) & 0.606 & 0.123 & 2.10 & Average predictions \\
044 & Residual from CF & 0.610 & 0.120 & 2.12 & CF + learned residual \\
045 & Long context (4s) & 0.606 & 0.120 & 2.10 & 4-second windows \\
046 & Calibrated output & 0.581 & 0.124 & 1.89 & Output calibration \\
047 & Spatial PE & 0.589 & 0.117 & 1.95 & REVE-inspired 4D PE \\
048 & NeuroTTT & 0.597 & 0.117 & 2.01 & Test-time SSL \\
049 & Adversarial & 0.608 & 0.117 & 2.12 & Gradient reversal \\
050 & L1 loss & 0.605 & 0.122 & 2.09 & MAE objective \\
051 & Per-channel heads & 0.609 & 0.122 & 2.14 & 12 separate decoders \\
052 & Cross-attention & 0.596 & 0.116 & 2.00 & Output queries \\
053 & Spectral loss & 0.604 & 0.121 & 2.04 & Frequency-domain MSE \\
054 & Pretrained (20 subj) & 0.589 & 0.117 & 1.94 & HBN encoder \\
\textbf{057} & \textbf{Unified pretrained} & 0.590 & 0.117 & 1.95 & 55-subject encoder \\
\textbf{078} & \textbf{Pruned subjects} & 0.610 & 0.120 & 2.11 & Drop subjects 1,2,4,7 \\
\bottomrule
\end{tabular}
\end{table}

Full 15-subject LOSO reference (CF baseline): $r = 0.645$,
95\% CI $= [0.563, 0.728]$, $\sigma = 0.144$.

%==========================================================================
\section{Key Findings and Implications}
\label{sec:findings}
%==========================================================================

\begin{enumerate}
  \item \textbf{The 3-subject benchmark was noise.} With CI width 0.47 and
    minimum permutation $p = 0.125$, the entire narrowband leaderboard
    (iter009--036) operated below measurement resolution. Future work must
    use 15-fold LOSO exclusively.

  \item \textbf{Pretraining quantity does not overcome architecture
    mismatch.} The unified encoder (55 subjects, val\_L1 $= 0.661$)
    improved reconstruction by 7\% over the 20-subject pilot, but
    fine-tuned performance gained only $+0.001\,r$. The cross-attention
    decoder is the bottleneck, not the encoder.

  \item \textbf{CF and deep models respond oppositely to data pruning.}
    Removing low-value subjects improves CF ($+0.013$) but hurts deep
    models ($-0.028$). Deep models extract useful temporal features even
    from subjects with atypical spatial mappings.

  \item \textbf{We are at 94\% of the cross-subject ceiling.} The CF
    baseline ($r = 0.645$) is within $+0.043\,r$ of the infinite-data
    asymptote ($r = 0.688$). The remaining $+0.10\,r$ to the noise ceiling
    ($0.791$) requires subject-specific adaptation.

  \item \textbf{25 diverse approaches produced zero improvement.} The
    robustness of the $r = 0.638$ ceiling across augmentation,
    meta-learning, domain adaptation, geometric methods, attention,
    data-centric pruning, regularization, ensembles, and pretraining
    suggests it is a genuine performance boundary for cross-subject
    prediction without target-subject calibration data.
\end{enumerate}

%==========================================================================
\section{Future Directions}
\label{sec:future}
%==========================================================================

The negative results from this phase point clearly toward two high-impact
directions:

\subsection{Subject-Specific Adaptation (Highest Priority)}

The $+0.10\,r$ gap between the cross-subject ceiling (0.688) and the
noise ceiling (0.791) is the largest remaining opportunity. Approaches:

\begin{itemize}
  \item \textbf{Few-shot calibration}: Use 30--60 seconds of
    target-subject data to adapt the spatial mapping via low-rank
    adaptation or affine transformation of the output layer.
  \item \textbf{Test-time channel weighting}: Estimate per-channel
    coupling strength from a brief calibration segment and mask out
    poorly-coupled channels.
  \item \textbf{Online adaptation}: Continuously update the model using
    recent predictions and self-supervised objectives.
\end{itemize}

\subsection{Direct Decoder Architecture}

The pretrained encoder failure (Section~\ref{sec:pretrain}) was caused by
the cross-attention decoder bottleneck. Future pretraining efforts should:

\begin{itemize}
  \item Use the MultiScaleConv+Transformer architecture end-to-end,
    initializing the convolutional layers from pretrained temporal features.
  \item Avoid cross-attention decoders that compress spatial information
    into a small number of query vectors.
  \item Consider direct feature transplant: freeze pretrained conv layers,
    train only the spatial mixing and output projection.
\end{itemize}

\subsection{Large-Scale Pretraining (TUH EEG Corpus)}

If pretraining is revisited, the TUH EEG Corpus ($\sim$21,000 hours,
30,000+ recordings) is the minimum viable dataset. The 55-subject unified
corpus was insufficient. Access requires institutional agreement and
$\sim$330\,GB of storage.

%==========================================================================
\section{Conclusion}
\label{sec:conclusion}
%==========================================================================

Iterations 055--079 represent the most comprehensive exploration of
training strategies for this task to date: 25 iterations spanning 7
technique categories, unified pretraining on 55 subjects across 5 datasets,
40+ research documents, and 7 papers analyzed. The unanimous negative
result---zero improvement over the 55K-parameter tiny deep baseline
($r = 0.638$)---is itself the most important finding. It establishes, with
high confidence, that the $r = 0.638$ ceiling is a genuine boundary for
cross-subject prediction on this data.

The statistical analysis revealed that the original 3-subject evaluation
was fundamentally flawed, retroactively invalidating the fine-grained
narrowband leaderboard. The data-centric analysis showed we are at 94\%
of the cross-subject ceiling, with the remaining improvement achievable
only through subject-specific adaptation. These findings redirect future
work away from architecture search and toward few-shot calibration and
large-scale pretraining---the only paths that can close the $+0.10\,r$ gap
to the noise ceiling.

\end{document}
